\section{Introduction}\label{sec:intro}

If a complex manifold is a product, its holomorphic tangent bundle is the
direct sum of the tangent bundles of the factors. The converse starts
with a specified holomorphic decomposition
\[
T_X=E_1\oplus E_2
\]
and asks whether the infinitesimal factors come from a product on the
universal cover. Compatibility with the given subbundles is part of the
question: an abstract product decomposition need not realize the
specified embeddings of $E_1$ and $E_2$ into $T_X$.
A subbundle is \emph{integrable} if its local holomorphic sections are
closed under the Lie bracket. When both summands are integrable,
holomorphic Frobenius gives simultaneous local product coordinates.
The problem is to continue these local products globally.

This question is closely related to the distinction between a local
product structure and a complete product geometry. For a complete
Riemannian metric, parallel complementary orthogonal distributions lead
to the de Rham decomposition \cite{deRham52}. A holomorphic splitting supplies no
complete metric making the distributions parallel. Compactness makes
every smooth metric complete, but it does not turn the given foliations
into parallel distributions. The missing step is therefore a
continuation theorem for the specified local products.

\subsection*{The compatible splitting problem}

Beauville formulated the compatible
universal-cover conjecture for compact K\"ahler manifolds with
integrable partial sums of the tangent summands
\cite[Section~2.3]{Beauville00}, and proved major positive cases, including
K\"ahler--Einstein manifolds and compact K\"ahler surfaces
\cite[Theorems~A and~C]{Beauville00}. In the K\"ahler--Einstein case,
Beauville refines an earlier splitting result of Yau \cite{Yau93}, as
he explains in his introduction. For two summands, the partial-sum
condition means precisely that both are integrable. For this formulation,
see also
\cite[Conjecture~1.4]{Horing26}.

The distinction between a product and a compatible product is already
visible in Beauville's example on $A\times\Pp^1$, where $A$ is an
abelian surface \cite[Section~2.2]{Beauville00}. Take independent
translation-invariant vector fields $U,V$ on $A$ and vector fields
$S,T$ on $\Pp^1$ with $[S,T]\ne0$, viewed as fields on the product.
The rank-two subbundle spanned by $U+S,V+T$ complements $T_{\Pp^1}$,
but is not integrable: its two generators have the nonzero vertical
bracket $[S,T]$. Thus even an underlying product need not realize a
specified tangent splitting.

Several later results establish global continuation under additional
geometric hypotheses. Druel treats projective manifolds whose tangent
bundle is a sum of line bundles, with the relevant partial sums
integrable, and obtains automatic integrability in the minimal
projective case \cite[Theorems~1 and~2]{Druel00}. Brunella, Pereira, and Touzet treat the case of a
line summand with integrable complement on a compact K\"ahler manifold
\cite[Theorem~1.1(i)]{BPT06}. Their work also gives uniformization
results when the tangent bundle is a sum of line bundles. Their proof
uses leafwise Poincar\'e metrics and transverse invariant metrics to
control continuation. Theorem~\ref{thm:main} below permits both summands
to have arbitrary positive rank.

In higher rank, Pereira and Touzet obtain a compatible Euclidean factor
when one involutive direct summand is Hermitian flat
\cite[Theorem~B]{PT13}. Druel, Pereira, Pym, and Touzet show that a
regular foliation with numerically flat tangent bundle---meaning that
the bundle and its dual are nef---is Hermitian flat and admits an
integrable complement, again producing a compatible Euclidean factor
\cite[Theorem~1.1]{DPPT24}.
A different continuation mechanism applies when a foliation has a compact
leaf. Druel, Pereira, Pym, and Touzet prove compatible
universal-cover splitting when one foliation has a compact leaf with
finite holonomy \cite[Theorem~1.4]{DPPT22}. Their argument uses global
Reeb stability and a flat Ehresmann connection. Our proof replaces
these additional flatness or leaf hypotheses by extension estimates
for maps adapted to the two distributions.

In the projective setting, integrability and product decomposition are
distinct issues. H\"oring proves automatic integrability for split
tangent bundles on non-uniruled projective manifolds
\cite[Theorem~1.3]{Horing07}. His proof combines the positivity theorem
for cotangent quotients of Campana and Peternell
\cite[Theorems~0.1 and~1.8]{CP11} with Demailly's integrability theorem
for holomorphic forms valued in the inverse of a pseudoeffective line
bundle \cite[Main Theorem]{Demailly02}. His recent work proves algebraic
integrability for tangent summands on normal $\mathbf Q$-factorial
projective klt varieties of Fano type; its product corollary, which does
not require $\mathbf Q$-factoriality, gives compatible products after
finite quasi-\'etale covers
\cite[Theorem~1.1 and Corollary~1.2]{Horing26}.
The companion \cite[Theorem~1.1]{OpenAIRCSplitting2026} proves that,
on a smooth connected rationally connected projective complex manifold,
both summands of every specified holomorphic splitting
$T_X=E_1\oplus E_2$ into positive-rank subbundles are integrable.
These results address
the geometry that forces a splitting to define foliations. In our main
theorem integrability is an explicit hypothesis, while the ambient
manifold is only assumed to be compact K\"ahler.

\begin{theorem}\label{thm:main}
Let $X$ be a compact connected K\"ahler manifold of complex dimension
$n\geq2$, and let
\[
T_X=E_1\oplus E_2
\]
be a specified decomposition into integrable holomorphic subbundles of
positive ranks $r_1,r_2$. If $\pi:\widetilde X\to X$ is the ordinary
universal covering map, there are connected simply connected complex
manifolds $Y_1,Y_2$ and a biholomorphism
$\Phi:\widetilde X\to Y_1\times Y_2$ such that
\[
\dd\Phi(\pi^*E_i)=\pr_i^*T_{Y_i}\qquad(i=1,2),
\]
where $\dd\pi$ identifies $T_{\widetilde X}$ with $\pi^*T_X$.
In particular, $\dim_\C Y_i=r_i$.
\end{theorem}

Theorem~\ref{thm:main} proves the two-summand form of Beauville's
compatible splitting conjecture. Its factors are the universal covers
of the intrinsic leaves through a common point. The leaves and factors
can be noncompact. The conclusion is about the ordinary universal cover;
it does not assert that a finite cover of $X$ is a product. If a summand
has rank zero, the corresponding factor is a point, so the positive-rank
formulation removes only the trivial case.

\begin{corollary}[Non-uniruled projective manifolds]\label{cor:non-uniruled}
Let $X$ be a smooth connected non-uniruled projective complex manifold,
and let $T_X=E_1\oplus E_2$ be a specified holomorphic splitting into
subbundles of positive ranks. Then both summands are integrable and the
ordinary universal cover has a compatible product decomposition as in
Theorem~\ref{thm:main}.
\end{corollary}

\begin{proof}
H\"oring's theorem \cite[Theorem~1.3]{Horing07} makes both summands
integrable. Since a smooth projective complex manifold is compact
K\"ahler, Theorem~\ref{thm:main} applies.
\end{proof}

\subsection*{The nef-and-big application}

For a smooth projective manifold, a line bundle is \emph{nef} if it has
nonnegative degree on every curve, and it is \emph{big} if its sections
have maximal asymptotic growth. When $K_X$ is nef and big, its
pluricanonical sections define a birational contraction $f:X\to Z$ to
the canonical model. We call $X^0=f^{-1}(Z_{\reg})$, the inverse image
of the smooth locus of $Z$, the \emph{regular canonical locus}.

Corollary~\ref{cor:non-uniruled} already gives the product conclusion
in this case. Our metric argument proves an additional statement:
the canonical K\"ahler--Einstein connection preserves both supplied
summands on $X^0$. This parallelism is the input for the alternative
continuation method in the appendices. We first record the product
conclusion.

\begin{corollary}\label{thm:nef-big}
Let $X$ be a smooth connected projective complex manifold of dimension
$n\geq2$ with nef and big canonical bundle $K_X$. Every specified
holomorphic splitting $T_X=E_1\oplus E_2$ into positive-rank subbundles
is integrable and induces a compatible product decomposition of the
ordinary universal cover as in Theorem~\ref{thm:main}.
\end{corollary}

When $K_X$ is ample, the negative K\"ahler--Einstein metric and
Beauville's theorem already give this conclusion. For nef and big $K_X$,
we use the canonical K\"ahler--Einstein approximation on a smooth
crepant resolution, as developed by Song \cite{Song14}.
The integrated squared second fundamental forms of the supplied summands
are $O(\varepsilon)$ along the approximation. Their limit vanishes on
the regular canonical locus, where the limiting metric makes both
summands parallel. The holomorphic Frobenius obstructions then vanish
on all of $X$. This supplies the precise hypothesis needed to apply
Theorem~\ref{thm:main}.
The approximate K\"ahler--Einstein method belongs to the stability
theory developed by Enoki \cite{Enoki88}; the passage from vanishing
second fundamental forms to a parallel splitting is used explicitly
by Guenancia \cite[Section~3.4, Step~4]{Guenancia16}.

\subsection*{The analytic continuation argument}

Local product coordinates alone do not show that a path in one foliation
can be transported all the way along a path in the other. We first solve
that extension problem for holomorphic maps, and then pass to continuous
paths. A holomorphic \emph{box} in two scalar parameters is a map $f$
whose $s$-derivative lies in $E_1$ and whose $t$-derivative lies in $E_2$.
In a simultaneous product chart it has the separated form
$(f_1(s),f_2(t))$. We do not impose an immersion or full-rank condition.

The first analytic step fills a Hartogs figure by a holomorphic box
(Proposition~\ref{prop:hartogs}). Ivashkovich's theorem gives a
meromorphic extension to the envelope of holomorphy
\cite{Ivashkovich1992,Ivashkovich2008}. The separate tangencies then
exclude positive-dimensional fibres of its graph. If a first target
coordinate varied along such a fibre, fixing that coordinate together
with $s$ would give an isolated level at a general fibre point. Nearby
ordinary graph points instead lie on one-dimensional levels, because
the same coordinate is independent of $t$. A compact local analytic
level would result, which is impossible in a coordinate ball. The
argument works for every summand rank and for boxes with deficient rank.

The second step crosses a smooth convex boundary at which the complex
tangent has nonzero components in both parameter directions
(Proposition~\ref{prop:mixed}). Bounded projections onto $E_1$ and
$E_2$ compare the normal and tangential speeds of the box. Stokes'
theorem for a K\"ahler form then bounds the areas of shrinking complex
tangent disks. Among many disjoint annuli, one has small area; a local
derivative rescaling argument gives a small image. Product coordinates
on that image build a Hartogs shell whose size exceeds the remaining
distance to the boundary. The first step fills it. Closedness of the
K\"ahler form enters the area estimate, while its metric need not make
the two summands orthogonal.

These extension results combine two holomorphic disk axes into a box
on their entire bidisk. A finite coordinate construction proves
holomorphic dependence on additional parameters, yielding full
polydisk boxes (Section~\ref{sec:boxes}). This is stronger than
continuation of individual germs: after transport, an edge still has
a holomorphic representative on its whole original parameter polydisk.
Consequently a finite rectangular grid combines arbitrary continuous
plaque paths, and a second grid realizes any continuous path as the
diagonal of a rectangle (Section~\ref{sec:global}). The rectangles
continue both the local product map and its inverse. Monodromy and
basepoint-matched covering lifts then give inverse biholomorphisms on
the simply connected covers.

\subsection*{An alternative through collapsed fibres}

The canonical contraction also permits a different continuation method.
We develop it in the appendices because it supplies a useful
neighborhood theorem for singular projective sets as well as another
proof of Corollary~\ref{thm:nef-big}. If $P$ is a compact simply connected
projective set in a smooth projective manifold with complementary integrable tangent
subbundles, and every two points of $P$ are joined by rational chains,
ordinary local crossing extends to a holomorphic map $c:U\times U\to M$
on a neighborhood of the whole set $P$
(Theorem~\ref{thm:crossing}).

The proof first fills products of rational curves, using meromorphic
Hartogs extension and the same complementary-tangency obstruction to
exceptional curves. It then uses proper stable-map spaces
\cite{FP97}, semialgebraic path lifting, and monodromy to continue
\emph{ambient} germs near the possibly singular or reducible set $P$.
For the canonical contraction, Hacon--McKernan's rational-chain theorem
and Takayama's local simple-connectedness theorem provide the required
fibre topology \cite{HM07,Takayama03}. Song's metric completion supplies
a continuous pseudodistance that collapses exactly those fibres.
Crossings are separately nonexpanding for this pseudodistance, and a
finite grid gives path rectangles even for paths contained in a
collapsed fibre. No length bound on those paths is required.

The main proof occupies Sections~\ref{sec:prelim}--\ref{sec:global}.
Section~\ref{sec:canonical} proves the metric estimate and the
nef-and-big application. Appendices~\ref{app:collapsed}--\ref{app:global}
give the alternative continuation mechanism in its dependency order.
All manifolds and varieties are complex. All covers and fundamental
groups are taken in the ordinary topology, and leaves carry their
intrinsic manifold topology.
